\begin{code}[hide]
module slides.Tsinghua-PL-seminar.sections.design where

open import Data.Nat using (ℕ ; _+_)
open import Data.Product using (_×_ ; _,_)
open import Data.Sum using (_⊎_ ; inj₁ ; inj₂)
open import Data.List using (List)
open import Data.Unit using (⊤ ; tt)
open import Level using (Level ; _⊔_ ; 0ℓ)
open import Relation.Binary using (Rel)
open import Relation.Binary.Definitions using (Decidable)
import Relation.Binary.PropositionalEquality as Eq
open import Function.Definitions using (Congruent ; Injective)
open import Algebra.Bundles using (Group)
open import Algebra.Morphism.Structures using (module GroupMorphisms)

open import Notations using (₁₊ ; ₂₊)
open import Word.Base using (wconcatmap)
import Word.Base as WB
open import Presentation.GroupLike using (Grouplike ; module Group-Lemmas)
open import Presentation.Construct.Base using ([_]ₗ ; [_]ᵣ)

postulate
  proof-elided : ∀ {ℓ} {A : Set ℓ} → A
  X Y A B C D  : Set
  a ℓ c d      : Level
  n k          : ℕ
\end{code}
%% Part 4 — the library, layer by layer (paper §4).
\PaperPart
\section{The library, layer by layer}

\begin{frame}{Five design decisions \InPaper}
\begin{columns}[c]
\begin{column}{0.60\textwidth}
\begin{enumerate}
  \setlength{\itemsep}{2pt}
  \item \textbf{Syntax is unquotiented}; everything out of it is a function
        on words --- normal forms, coset tables, sections are structurally
        recursive programs the evaluator runs.
  \item \textbf{Specifications are stdlib bundles}: a normal form \emph{is}
        an \aid{Injection}, \aid{RightInverse} or \aid{Bijection} out of the
        word setoid. No new theory to maintain.
  \item \textbf{Certificates are finite}: a verified coset table is finite
        data plus five equations; obligations are case splits whose cases
        are closed goals.
  \item \textbf{Indices are chosen for the unifier}: wire counts built from
        successors only.
  \item \textbf{Constructions come in pairs}: a syntactic theorem (normal
        forms combine) and a semantic one (the constructed relation set
        presents the constructed group).
\end{enumerate}
\end{column}
\begin{column}{0.37\textwidth}
\centering
\begin{tikzpicture}[x=10mm,y=7.6mm]
  \node[ly=paperC!25,minimum width=50mm] at (2.5,7) {Examples: $S_n$, Pauli, Clifford+$T$, \ldots};
  \node[ly=soft,minimum width=24mm] at (1.2,6) {Constructions\\[-2pt]{\tiny $\times,\rtimes,\ast_M$, ext.}};
  \node[ly=soft,minimum width=24mm] at (3.8,6) {Circuits\\[-2pt]{\tiny \aid{Lift-Relation}}};
  \node[ly=soft,minimum width=50mm] at (2.5,5) {Normalization: NF bundles, R--S, cosets, towers};
  \node[ly=soft,minimum width=50mm] at (2.5,4) {Presentations: $\_{===}\_$ and its congruence $\_{\approx}\_$};
  \node[ly=soft,minimum width=50mm] at (2.5,3) {Words: the free monoid \aid{Word X}};
  \node[ly=gray!10,minimum width=50mm] at (2.5,2) {\texttt{Notations} \quad\textbar\quad \texttt{ForStdlib} (semantics)};
  \node[ly=gray!5,draw=gray,minimum width=50mm] at (2.5,1) {Agda standard library 2.4};
\end{tikzpicture}
\end{column}
\end{columns}
\end{frame}

\begin{frame}{Words: a syntax tree, not a list \InPaper}
\begin{columns}[c]
\begin{column}{0.56\textwidth}
\begin{code}
data Word (X : Set) : Set where
  [_]ʷ : X → Word X
  ε    : Word X
  _•_  : Word X → Word X → Word X
\end{code}
\begin{code}[hide]
infixr 7 _•_
infixl 9 _ʷ _ᵗ
infix 4 _≈_ _===_ _≈ₙ_ _≈₂_ _~_ _===₁_ _===₂_


WRel : Set → Set₁
WRel X = Rel (Word X) 0ℓ

postulate
  w v u w' v' : Word X
  _===_       : WRel X
\end{code}
A word remembers its bracketing: terms mirror pen-and-paper calculations,
rules apply to any subterm, and associativity is a \alert{congruence
generator}, discharged by a solver.

\vspace{4pt}
Generator substitutions extend to words (``the unique monoid homomorphism''):
\begin{code}
_ʷ : (X → Word Y) → (Word X → Word Y)
_ʷ = wconcatmap
\end{code}
and a \alert{coset action} extends to words by threading the coset:
\begin{code}
_ᵗ : (C → Y → Word X × C) → (C → Word Y → Word X × C)
\end{code}
\begin{code}[hide]
_ᵗ = WB._ᵗ
\end{code}
\end{column}
\begin{column}{0.41\textwidth}
\centering
\begin{tikzpicture}[level distance=9mm,sibling distance=14mm,
  every node/.style={font=\small},
  op/.style={circle,draw=accent,fill=soft,inner sep=1pt,minimum size=5mm},
  lf/.style={rectangle,draw=paperC,fill=softpaper,inner sep=2pt,rounded corners=1pt}]
  \node[op] {$\bullet$}
    child {node[op] {$\bullet$}
      child {node[lf] {$H$}}
      child {node[lf] {$S$}}}
    child {node[op] {$\bullet$}
      child {node[lf] {$H$}}
      child {node[lf,draw=gray,fill=white] {$\varepsilon$}}};
\end{tikzpicture}

{\scriptsize $(H \aidop{•} S) \aidop{•} (H \aidop{•} \varepsilon)$ --- a tree;\\
\aid{by-assoc} flattens both sides to lists and compares by \aid{refl}}

\vspace{8pt}
\begin{tikzpicture}[>={Stealth[length=1.8mm]}]
  \node[draw=accent,fill=soft,rounded corners=2pt,font=\scriptsize] (c) {$c$};
  \node[draw=accent,fill=soft,rounded corners=2pt,font=\scriptsize,right=12mm of c] (c1) {$c_1$};
  \node[draw=accent,fill=soft,rounded corners=2pt,font=\scriptsize,right=12mm of c1] (c2) {$c_2$};
  \draw[->] (c) -- node[above,font=\scriptsize]{$y_1$} node[below,font=\scriptsize]{emit $w_1$} (c1);
  \draw[->] (c1) -- node[above,font=\scriptsize]{$y_2$} node[below,font=\scriptsize]{emit $w_2$} (c2);
\end{tikzpicture}

{\scriptsize $(h\,\aid{ᵗ})\;c\;(y_1\aidop{•}y_2) = (w_1\aidop{•}w_2,\;c_2)$}
\end{column}
\end{columns}
\end{frame}

\begin{frame}{Presented monoids as inductive congruences \InPaper}
\begin{columns}[c]
\begin{column}{0.52\textwidth}
One convention everywhere: \aid{\_===\_} = the raw axioms $\Gamma$;
\aid{\_≈\_} = the monoid congruence they generate:
\begin{code}
data _≈_ : WRel X where
  refl  : w ≈ w
  sym   : w ≈ v → v ≈ w
  trans : w ≈ v → v ≈ u → w ≈ u
  cong  : w ≈ w' → v ≈ v' → w • v ≈ w' • v'
  assoc      : (w • v) • u ≈ w • (v • u)
  left-unit  : ε • w ≈ w
  right-unit : w • ε ≈ w
  axiom      : w === v → w ≈ v
\end{code}
\end{column}
\begin{column}{0.45\textwidth}
\begin{itemize}
  \setlength{\itemsep}{4pt}
  \item $\approx$ \emph{is} the word problem of $\langle X\mid\Gamma\rangle$.
  \item A derivation is a \alert{first-class object}: compute with it,
        induct on it, translate it (congruence-lifting lemmas are inductions
        on this type).
  \item Packaged as the \emph{word setoid} and the presented monoid
        \aid{•-ε-monoid}; with a \aid{Grouplike} witness (a left inverse per
        generator) the group \aid{•-ε-group}.
\end{itemize}
\vspace{2pt}
{\small Contrast: mathlib's \aid{PresentedGroup} (quotient), cubical Agda (HIT).}
\end{column}
\end{columns}
\end{frame}

\begin{frame}{What ``$\Gamma$ presents $G$'' means \InPaper}
\begin{code}
record _IsPresentationOf_ (_===_ : WRel X) (G : Group a ℓ) : Set (a ⊔ ℓ) where
  field gl : Grouplike _===_
  module GL = Group-Lemmas _===_ gl
  open GroupMorphisms (Group.rawGroup GL.•-ε-group) (Group.rawGroup G)
  field ⟦_⟧ : Word X → Group.Carrier G
        iso : IsGroupIsomorphism ⟦_⟧
\end{code}
\begin{code}[hide]
postulate
  |NF|       : Set
  nf         : Word X → |NF|
  _≈ₙ_       : |NF| → |NF| → Set
  NormalForm : Set
  inv-nf     : |NF| → Word X
  Sem        : Set
  ⟦_⟧        : Word X → Sem
  _≈₂_       : Sem → Sem → Set
\end{code}
\begin{columns}[c]
\begin{column}{0.54\textwidth}
\small
\begin{itemize}
  \item weaker records: \aid{\_IsSubPresentationOf\_} (monomorphism), monoid versions;
  \item upgrade lemmas: surjective sub-presentation $\Rightarrow$ presentation;
        grouplike monoid presentation $\Rightarrow$ group presentation;
  \item proofs split into soundness, injectivity (by normalization),
        surjectivity; the records assemble the isomorphism.
\end{itemize}
\end{column}
\begin{column}{0.43\textwidth}
\begin{keybox}[title=Monoid vs group: a real distinction]
\small Words have no formal inverses, so \emph{every} presentation is at
bottom a monoid presentation. $\langle\, T \mid T^{0} = \varepsilon\,\rangle$
read as a group presentation is $\mathbb{Z}$; as a \alert{monoid} it is
$(\mathbb{N},+,0)$ --- and \aid{Grouplike} is exactly what fails.
\end{keybox}
\end{column}
\end{columns}
\end{frame}

\begin{frame}{Normal forms are standard-library bundles \InPaper}
\begin{columns}[c]
\begin{column}{0.50\textwidth}
\small From the \alert{word setoid} to a carrier $\mathit{NF}$, on the nose:\\[4pt]
\begin{tabular}{@{\quad}l@{\;\;=\;\;}l@{}}
\aid{NormalFormInjective} & \aid{Injection} \\
\aid{NormalForm}          & \aid{RightInverse} \\
\aid{BijectiveNormalForm} & \aid{Bijection}
\end{tabular}
\vspace{6pt}


\aid{Injection}: $w\approx v \iff \aid{nf}\,w \approx_{\mathit{NF}} \aid{nf}\,v$, and with it
\begin{code}
by-equal-nf : nf w ≈ₙ nf v → w ≈ v
≈-dec : Decidable _≈ₙ_ → Decidable _≈_
\end{code}
\begin{code}[hide]
by-equal-nf = proof-elided
≈-dec = proof-elided
\end{code}
{\small prove word equations by \alert{computing}; a decidable carrier
\alert{decides the word problem}.}
\end{column}
\begin{column}{0.47\textwidth}
\centering
\begin{tikzpicture}[>={Stealth[length=2mm]}]
  \node[draw=accent,fill=soft,ellipse,minimum width=26mm,minimum height=34mm] (W) at (0,0) {};
  \node[draw=paperC,fill=softpaper,ellipse,minimum width=18mm,minimum height=28mm] (N) at (4,0) {};
  \node[font=\scriptsize\bfseries,accent] at (0,1.95) {word setoid};
  \node[font=\scriptsize\bfseries,paperC] at (4,1.95) {$\mathit{NF}$};
  \foreach \y in {0.9,0.35,-0.2} \node[circle,fill=accent,inner sep=1pt] at (-0.4,\y) {};
  \draw[accent!60,rounded corners] (-0.8,-0.45) rectangle (0.1,1.15);
  \node[circle,fill=paperC,inner sep=1.3pt] (u) at (4,0.45) {};
  \draw[->,thick] (0.15,0.45) to[bend left=10] node[above,font=\scriptsize\agdafont]{nf} (u);
  \node[circle,fill=newC,inner sep=1.3pt] (r) at (-0.4,-1.0) {};
  \draw[->,thick,newC] (u) to[bend left=25] node[below right,font=\scriptsize\agdafont]{inv-nf} (r);
  \node[font=\tiny,newC,align=center] at (-0.4,-1.35) {canonical\\representative};
  \node[circle,fill=gray,inner sep=1pt] at (4.3,-1.0) {};
  \node[font=\tiny,gray,align=center] at (4.3,-1.35) {junk allowed};
\end{tikzpicture}

{\small \aid{RightInverse} adds the \alert{section} \aid{inv-nf} with
$\aid{inv-nf}\circ\aid{nf}\approx\mathrm{id}$. Exactness
$\aid{nf}\circ\aid{inv-nf}\approx\mathrm{id}$ is kept separate.}
\end{column}
\end{columns}
\end{frame}

\begin{frame}{Where constructivity bites: the section is data \InPaper}
\begin{columns}[c]
\begin{column}{0.55\textwidth}
\begin{itemize}
  \setlength{\itemsep}{5pt}
  \item Classically, an injective \aid{nf} has a section on its image for
        free: \emph{choose} a preimage.
  \item Constructively, choosing is not free: an \aid{Injection} gives
        \alert{no algorithm} producing a word that realizes a normal form.
  \item That algorithm is exactly the extra field of \aid{RightInverse} ---
        and it is what completeness consumes.
  \item So the two are \alert{different strengths} in Agda, though
        classically equal. Use the weak one for word-problem reasoning, the
        strong one for completeness and presentations.
  \item A fourth, crude witness, \aid{WeakNormalForm}: an injective
        \emph{invariant}, not even congruent.
\end{itemize}
\end{column}
\begin{column}{0.42\textwidth}
\centering
\begin{tikzpicture}[x=12mm,y=10mm]
  \node[w=gray] (weak) at (0,0) {\aid{WeakNormalForm}};
  \node[w=accent] (inj) at (0,1.3) {\aid{NormalFormInjective}};
  \node[w=paperC] (ri) at (0,2.6) {\aid{NormalForm}\\[-2pt]{\scriptsize + section (data)}};
  \node[w=newC] (bij) at (0,3.9) {\aid{BijectiveNormalForm}\\[-2pt]{\scriptsize surjectivity as a property}};
  \draw[-{Stealth},thick] (ri) -- (inj);
  \draw[-{Stealth},thick] (inj) -- (weak);
  \draw[-{Stealth},thick] (bij) -- (ri);
  \node[font=\scriptsize,anchor=west,align=left] at (1.55,1.95) {strictly stronger\\as \emph{data}};
\end{tikzpicture}
\end{column}
\end{columns}
\end{frame}

\begin{frame}{Completeness, once and for all \InPaper}
A \aid{NormalForm} is \alert{unique} for a semantics \aid{⟦\_⟧} into any setoid
when canonical representatives with equal denotations are equal:
\begin{code}
record UniqueNormalForm (normalForm : NormalForm) : Set (c ⊔ d) where
  field unique : ∀ {u v : |NF|} → ⟦ inv-nf u ⟧ ≈₂ ⟦ inv-nf v ⟧ → u ≈ₙ v
by-normalization : {normalForm : NormalForm} → UniqueNormalForm normalForm →
  Congruent _≈_ _≈₂_ ⟦_⟧ → Injective _≈_ _≈₂_ ⟦_⟧
\end{code}
\begin{code}[hide]
by-normalization = proof-elided
\end{code}
\begin{columns}[c]
\begin{column}{0.55\textwidth}
\small
\textbf{Soundness + unique NF $\Rightarrow$ completeness.} Five lines of proof;
its value: every case study only ever proves \aid{unique} --- a statement about
\alert{normal forms}, i.e.\ concrete finite data, never about arbitrary words.
\end{column}
\begin{column}{0.42\textwidth}
\small
\begin{itemize}
  \item \aid{by-normalization-wsurj}: surjectivity on canonical forms suffices;
  \item \aid{by-completeness}: the converse, given an exact section --- how
        uniqueness lifts through products.
\end{itemize}
\end{column}
\end{columns}
\end{frame}

\begin{frame}{Inductively defined circuits: indices for the unifier \InPaper}
\begin{columns}[c]
\begin{column}{0.52\textwidth}
\aid{\_+\_} recurses on its \emph{first} argument:
$2 + n \;\leadsto\; \mathsf{suc}\,(\mathsf{suc}\,n)$, while $n + 2$ is stuck.

\vspace{4pt}
So the indices are \alert{left-biased} too:
\begin{itemize}
  \item \aid{gate₂ : Gate 2 → Gen (₂₊ n)}, not \aid{Gen k} with $k\ge 2$;
  \item $k$-fold shift: \aid{Gen n → Gen (k + n)}.
\end{itemize}
Every index equation the coverage checker meets is solved by
\alert{injectivity of \aid{₁₊}}, never by arithmetic. That is why coset tables
are written by direct pattern matching.

\vspace{4pt}
\small \aid{lemma-cong↑} lifts the \emph{entire} congruence one wire up, by
induction over the eight constructors of $\approx$.
\end{column}
\begin{column}{0.45\textwidth}
\centering
\begin{tikzpicture}[
  ok/.style={draw=paperC,fill=softpaper,rounded corners=2pt,font=\small,align=left,text width=50mm,inner sep=4pt},
  bad/.style={draw=newC,fill=softnew,rounded corners=2pt,font=\small,align=left,text width=50mm,inner sep=4pt}]
  \node[ok] (a) {\textbf{left-biased} \hfill\checkmark\\[2pt]
   \aid{Gen (₂₊ n)} $\doteq$ \aid{Gen (₁₊ m)}\\
   $\leadsto\ m := \aid{₁₊}\,n$ \quad by unification};
  \node[bad,below=6mm of a] (b) {\textbf{right-biased} \hfill$\times$\\[2pt]
   \aid{Gen (n + 2)} $\doteq$ \aid{Gen (₁₊ m)}\\
   $\leadsto\ n + 2 \doteq \mathsf{suc}\,m$ \quad stuck:\\ a Diophantine equation};
\end{tikzpicture}
\end{column}
\end{columns}
\end{frame}

\begin{frame}{The Reidemeister--Schreier engine \InPaper}
\begin{code}
module SingleLevel (Γ : WRel X) (Δ : WRel Y) (C : Set) (I : C)
  (f : X → Word Y) (h : C → Y → Word X × C) ([_] : C → Word Y)
\end{code}
\begin{code}[hide]
  where

postulate
  I      : C
  f      : X → Word Y
  h      : C → Y → Word X × C
  [_]    : C → Word Y
  _~_    : Word X × C → Word X × C → Set
  _===₁_ : WRel X
  _===₂_ : WRel Y
\end{code}
\vspace{-2pt}
{\small $\Gamma$ presents a subgroup $H$ of the group $G$ presented by $\Delta$;
$C$ = cosets, $I$ = identity coset, $f$ = embedding of generators,
$[\_]$ = Schreier section, $h$ = the coset table.}
\begin{code}
module Transfer
  (h=⁻¹f-gen : ∀ (x : X) → ([ x ]ʷ , I) ~ ((h ᵗ) I (f x)))                         -- 1
  (h-wd-ax : ∀ (c : C){u t : Word Y} → u ===₂ t → ((h ᵗ) c u) ~ ((h ᵗ) c t))       -- 2
  (f-wd-ax : ∀ {w v} → w ===₁ v → (f ʷ) w ≈₂ (f ʷ) v)                              -- 3
  ([I]≈ε : [ I ] ≈₂ ε)                                                             -- 4
  (h=ract : ∀ c b → let (b' , c') = h c b in [ c ] • [ b ]ʷ ≈₂ (f ʷ) b' • [ c' ])  -- 5
\end{code}
\begin{code}[hide]
  where
\end{code}
\vspace{-4pt}
{\small (1) the table inverts the embedding at $I$;\; (2) it respects $\Delta$'s
axioms, coset by coset;\; (3) $f$ respects $\Gamma$'s axioms;\; (4) $[I]\approx\varepsilon$;\;
(5) the soundness law, as \aid{ract-sound}.}
\end{frame}

\begin{frame}{What the engine gives back \InPaper}
\begin{columns}[c]
\begin{column}{0.52\textwidth}
From the five hypotheses, \aid{Transfer} derives well-definedness of
$\aid{nf} = (h\,\aid{ᵗ})\,I$, the factorization of whole words, injectivity,
and the payoff --- the \alert{normal-form transport}:
\[
  \aid{nf-transp} : \aid{NormalForm}\;\Gamma\;\mathit{NF}
    \to \aid{NormalForm}\;\Delta\;(\mathit{NF}\times C)
\]
a normal form for the subgroup yields one for the group, \alert{with one
more coset digit}.

\vspace{4pt}
\small \aid{PackedCosetTable} bundles data + hypotheses into one value;
\aid{CosetTower} iterates up an $\mathbb{N}$-indexed family:\\[2pt]
\quad\aid{tower-carrier NF₀ (k+1) = tower-carrier NF₀ k × Cₖ}
\end{column}
\begin{column}{0.45\textwidth}
\centering
\begin{tikzpicture}[x=9mm,y=9mm,>={Stealth[length=1.8mm]},
  lv/.style={draw=accent,fill=soft,rounded corners=2pt,font=\scriptsize,minimum width=26mm,minimum height=6mm},
  dg/.style={draw=newC,fill=softnew,rounded corners=2pt,font=\scriptsize,minimum width=8mm,minimum height=6mm}]
  \node[lv] (l0) at (0,0) {$\mathit{NF}_0$ (base)};
  \node[lv] (l1) at (0,1.2) {level 1};
  \node[lv] (l2) at (0,2.4) {level 2};
  \node[lv] (l3) at (0,3.6) {level $k$};
  \node[dg] (d1) at (2.6,1.2) {$C_1$};
  \node[dg] (d2) at (2.6,2.4) {$C_2$};
  \node[dg] (d3) at (2.6,3.6) {$C_k$};
  \draw[->,thick,paperC] (l0) -- node[left,font=\tiny]{table 1} (l1);
  \draw[->,thick,paperC] (l1) -- node[left,font=\tiny]{table 2} (l2);
  \draw[->,thick,paperC,dashed] (l2) -- (l3);
  \foreach \i in {1,2,3} \draw[-,gray] (l\i) -- (d\i);
  \node[font=\scriptsize,align=center] at (1.3,-1.0) {carrier $\mathit{NF}_0\times C_1\times C_2\times\dots\times C_k$};
\end{tikzpicture}
\end{column}
\end{columns}
\end{frame}

